\documentclass[../DoAn.tex]{subfiles}
\begin{document}

\section{Mở đầu chương}
Sau khi hoàn thành việc huấn luyện mô hình TinyLLaMA sử dụng phương pháp LoRA trên tập dữ liệu câu hỏi – trả lời trích xuất từ email, chương này sẽ tiến hành đánh giá hiệu quả của hệ thống chatbot. Các nội dung chính bao gồm: đánh giá pipeline huấn luyện, phân tích hiệu suất và tài nguyên sử dụng, đánh giá chất lượng phản hồi, cùng những điểm mạnh và hạn chế còn tồn tại.

\section{Đánh giá pipeline huấn luyện mô hình}

\subsection{Ưu điểm nổi bật}

\textbf{1. Pipeline đầy đủ và hợp lý} \\ 
Quy trình huấn luyện được xây dựng theo cấu trúc rõ ràng, bao gồm:
\begin{itemize}
    \item Load dữ liệu từ định dạng JSONL.
    \item Format lại prompt theo chuẩn trò chuyện người – máy.
    \item Huấn luyện mô hình TinyLLaMA với kỹ thuật LoRA.
    \item Sinh phản hồi từ mô hình đã fine-tune.
\end{itemize}

\textbf{2. Cấu hình LoRA phù hợp với mô hình nhỏ} \\ 
Sử dụng \texttt{r=16}, \texttt{lora\_alpha=32}, \texttt{lora\_dropout=0.05} – các thông số cân bằng giữa tốc độ và khả năng học. Target modules gồm \texttt{q\_proj}, \texttt{v\_proj}, phù hợp với kiến trúc LLaMA. Ngoài ra, sử dụng \texttt{load\_in\_8bit=True} giúp giảm tới 50\% bộ nhớ GPU.

\textbf{3. Tối ưu bộ nhớ và tốc độ} \\ 
Dùng batch size nhỏ kết hợp với \texttt{gradient\_accumulation} giúp tránh hết VRAM khi chạy trên GPU T4 hoặc Colab. Với tập dữ liệu nhỏ, thời gian huấn luyện chỉ mất khoảng 30–60 phút cho 3 epoch.

\subsection{Hạn chế và các vấn đề cần cải thiện}

\textbf{1. Thiếu validation và đánh giá định lượng} \\ 
\begin{itemize}
    \item Dataset hiện tại chưa tách tập kiểm thử (\texttt{train\_test\_split}).
    \item Chưa áp dụng các chỉ số như \texttt{eval\_loss}, perplexity, BLEU,... để đo hiệu quả.
    \item Chưa có cấu hình lưu mô hình tốt nhất hay kiểm soát đánh giá theo từng bước.
\end{itemize}

\textbf{2. Chưa kiểm soát chất lượng dữ liệu đầu vào} \\ 
Một số email sau xử lý vẫn còn nhiễu hoặc chứa thông tin dư thừa. Không kiểm tra độ dài trung bình của câu dẫn đến padding dư thừa và token không cần thiết.

\textbf{3. Prompt chưa đủ tối ưu} \\ 
Hiện tại đang sử dụng định dạng đơn giản:
\begin{adjustwidth}{1em}{1em}
\begin{lstlisting}[language=none]
### User:
[instruction]

### Assistant:
[output]
\end{lstlisting}
\end{adjustwidth}
Cần cải thiện bằng cách dùng prompt hướng dẫn rõ hơn, chẳng hạn như:
\begin{adjustwidth}{1em}{1em}
\begin{lstlisting}[language=none]
<|system|>
Bạn là một trợ lý học vụ thông minh.

<|user|>
[instruction]

<|assistant|>
[output]<|end|>
\end{lstlisting}
\end{adjustwidth}

\textbf{4. Cấu hình sinh phản hồi chưa tối ưu} \\ 
Cần thêm các tham số như \texttt{top\_k}, \texttt{top\_p}, \texttt{temperature}, \texttt{repetition\_penalty} để cải thiện sự đa dạng và kiểm soát tính lặp lại.

\section{Đánh giá chất lượng phản hồi}

\subsection{Ưu điểm}
\begin{itemize}
    \item Ngôn ngữ phản hồi rõ ràng, đúng phong cách học thuật.
    \item Trả lời tốt các câu hỏi phổ biến đã xuất hiện trong dữ liệu huấn luyện.
    \item Phản hồi nhanh, phù hợp chạy trên máy cá nhân hoặc Google Colab.
\end{itemize}

\subsection{Nhược điểm}
\begin{itemize}
    \item Mô hình dễ trả lời sai khi gặp câu hỏi chưa từng xuất hiện.
    \item Các câu hỏi quá ngắn hoặc thiếu ngữ cảnh thường bị trả lời lệch chủ đề.
    \item Không có khả năng tra cứu thời gian thực (ví dụ: lịch học mới).
    \item Thiếu khả năng suy luận đa bước (multi-hop reasoning).
\end{itemize}

\section{Gợi ý cải tiến}

\begin{itemize}
    \item Tách tập validation để có thể đo \texttt{eval\_loss}, \texttt{perplexity}
    \item áp dụng \texttt{load\_best\_model\_at\_end}.
    \item Thêm \texttt{EarlyStoppingCallback} để tránh overfitting.
    \item Cải tiến prompt theo phong cách hệ thống – người dùng – trợ lý.
    \item Sử dụng cấu hình sinh nâng cao:
    \begin{adjustwidth}{1em}{1em}
    \begin{lstlisting}[language=none,breaklines=true,basicstyle=\ttfamily\footnotesize]
generation_config = {
  "max_new_tokens": 512,
  "do_sample": True,
  "temperature": 0.7,
  "top_p": 0.9,
  "top_k": 50,
  "repetition_penalty": 1.1
}
\end{lstlisting}
    \end{adjustwidth}
\end{itemize}
\section{Kết luận chương}

Mô hình TinyLLaMA fine-tune bằng LoRA là một lựa chọn hợp lý cho các tác vụ chatbot học vụ nhẹ, đặc biệt trong bối cảnh tài nguyên giới hạn. Pipeline huấn luyện rõ ràng, khả năng mở rộng tốt và phù hợp cho các bài toán Q\&A đơn giản. Tuy nhiên, để đạt hiệu quả cao hơn, hệ thống cần được cải thiện về chất lượng dữ liệu, đánh giá định lượng và prompt thiết kế.

\textbf{Phù hợp cho:} chatbot Q\&A chuyên ngành, ứng dụng nhỏ. \\
\textbf{Không phù hợp cho:} tác vụ suy luận phức tạp, môi trường production lớn.

\end{document}